\chapter{The Weil Conjectures from the Perverse Package}
\label{chap:weil-assembly}

Let \(X_0\) be smooth, projective, and geometrically connected of dimension
\(d\) over \(k_0=\Fq\), and put
\[
  X=X_0\times_{\Fq}\Fbar,\qquad
  a_0:X_0\longrightarrow\Spec\Fq.
\]
Fix an auxiliary prime \(\ell\ne\operatorname{char}\Fq\), and let \(F\) be
geometric Frobenius.

\section{The one-line perverse proof of purity}

\begin{theorem}[Terminal perverse calculation]
\label{thm:terminal-perverse-calculation}
Every eigenvalue \(\alpha\) of \(F\) on
\[
  V_{i,\ell}=H^i_{\et}(X,\Ql)
\]
has \(\iota\)-absolute value \(q^{i/2}\), for every field isomorphism
\(\iota:\Qbarl\simeq\mathbf C\).
\end{theorem}

\begin{proof}
The constant sheaf is punctually pure of weight \(0\).  Since \(X_0\) is
smooth of dimension \(d\),
\[
  \IC_{X_0}=\Qbarl[d]
\]
is perverse and pure of weight \(d\), by
\cref{eq:IC-smooth,eq:constant-perverse-weight}.  The structural morphism is
proper, so \cref{thm:proper-direct-image-purity} makes
\[
  Ra_{0*}\Qbarl[d]
\]
pure of weight \(d\).  On the point, perverse and ordinary cohomology agree:
\[
  \pH^r(Ra_{0*}\Qbarl[d])
  =H^r(Ra_{0*}\Qbarl[d])
  =H^{r+d}(X,\Qbarl).
\]
The left side is pure of weight \(d+r\).  Taking \(r=i-d\) gives weight
\(i\), which means \(|\iota(\alpha)|=q^{i/2}\).

The argument in the weight chapters was carried out for an arbitrary
\(\iota\); repeating it for every \(\iota\) gives the universal-realization
assertion.  No rationality or algebraicity of the individual factor has yet
been assumed.
\end{proof}

\begin{perversebox}
The decomposition theorem is compatible with this calculation but not needed
for it.  Proper purity is the essential input.  Logically, decomposition was
proved only after proper purity in \cref{chap:decomposition}.
\end{perversebox}

\section{Trace formula and rationality}

Finiteness gives \(V_{i,\ell}=0\) outside \(0\leq i\leq2d\), and all these
spaces are finite dimensional.  Properness identifies ordinary and
compactly supported cohomology.  The trace formula applied to \(F^n\) gives
\[
  \#X_0(\mathbf F_{q^n})
  =
  \sum_{i=0}^{2d}(-1)^i
  \Tr(F^n\mid V_{i,\ell}).
\tag{17.1}\label{eq:trace-X}
\]
Put
\[
  P_{i,\ell}(T)=\det(1-TF\mid V_{i,\ell})
  \in1+T\Ql[T].
\tag{17.2}\label{eq:Pi-definition}
\]

\begin{theorem}[Cohomological rationality]
\label{thm:cohomological-rationality}
For every \(\ell\ne\operatorname{char}\Fq\),
\[
  Z(X_0,T)
  =
  \prod_{i=0}^{2d}P_{i,\ell}(T)^{(-1)^{i+1}}
  \quad\text{in }\Ql[[T]].
\tag{17.3}\label{eq:cohomological-zeta}
\]
In particular \(Z(X_0,T)\in\mathbf Q(T)\).
\end{theorem}

\begin{proof}
Insert \eqref{eq:trace-X} into
\[
 Z(X_0,T)=
 \exp\!\left(\sum_{n\geq1}
 \#X_0(\mathbf F_{q^n})\frac{T^n}{n}\right).
\]
The alternating determinant--trace identity
\cref{cor:alternating-det-trace} changes the result into
\eqref{eq:cohomological-zeta}.  Thus \(Z\) is rational over \(\Ql\).

The point-counting definition lies in \(\mathbf Q[[T]]\), and the closed-point
Euler product \cref{prop:euler-product} lies in \(\mathbf Z[[T]]\).
Hankel descent \cref{thm:hankel-descent}, applied to
\(\mathbf Q\subset\Ql\), yields \(Z\in\mathbf Q(T)\).
\end{proof}

\begin{warning}
At this point only the alternating product is rational.  It would be false in
general to infer that each \(P_{i,\ell}\) is rational: numerator and
denominator factors can cancel.  Universal purity will exclude such
cancellation by separating their absolute values.
\end{warning}

\section{Endpoint factors}

\begin{proposition}[Degree zero and top degree]
\label{prop:endpoint-factors}
One has
\[
  P_{0,\ell}(T)=1-T,\qquad
  P_{2d,\ell}(T)=1-q^dT.
\tag{17.4}\label{eq:endpoint-factors}
\]
\end{proposition}

\begin{proof}
Geometric connectedness gives
\(H^0(X,\Ql)=\Ql\), with Frobenius acting trivially.  Poincare duality
identifies \(H^{2d}(X,\Ql)\) with \(\Ql(-d)\), on which geometric Frobenius
acts by \(q^d\).  These are exactly the two determinants.
\end{proof}

\section{Duality and reciprocal eigenvalues}

Poincare duality gives perfect \(F\)-equivariant pairings
\[
  V_{i,\ell}\times V_{2d-i,\ell}
  \longrightarrow\Ql(-d)
\tag{17.5}\label{eq:poincare-X}
\]
and therefore
\[
  \langle Fx,Fy\rangle=q^d\langle x,y\rangle.
\tag{17.6}\label{eq:poincare-similitude}
\]
If \(\alpha_1,\ldots,\alpha_{b_i}\) are the eigenvalues on \(V_{i,\ell}\),
then those on \(V_{2d-i,\ell}\) are
\[
  \frac{q^d}{\alpha_1},\ldots,\frac{q^d}{\alpha_{b_i}}.
\tag{17.7}\label{eq:reciprocal-eigenvalues}
\]
In particular \(b_i=b_{2d-i}\).

\begin{proposition}[Purity in every complex realization]
\label{prop:all-realization-purity}
Every eigenvalue \(\alpha\in\Qbarl^\times\) of geometric Frobenius on
\(V_{i,\ell}\) satisfies
\[
  |\iota(\alpha)|=q^{i/2}
\tag{17.8}\label{eq:smooth-proper-purity}
\]
for every isomorphism \(\iota:\Qbarl\simeq\mathbf C\).
\end{proposition}

\begin{proof}
This is \cref{thm:terminal-perverse-calculation}.  The word “eigenvalue” here
means a root of the characteristic polynomial of \(F\); the roots of
\(P_{i,\ell}(T)\) as a polynomial in \(T\) are their reciprocals.
\end{proof}

\section{The functional equation}

Let
\[
  \chi_\ell=\sum_{i=0}^{2d}(-1)^i\dim_{\Ql}V_{i,\ell}.
\]

\begin{theorem}[Functional equation]
\label{thm:zeta-functional-equation}
There exists \(\varepsilon_\ell\in\{1,-1\}\) such that
\[
  Z(X_0,q^{-d}T^{-1})
  =
  \varepsilon_\ell q^{d\chi_\ell/2}
  T^{\chi_\ell}Z(X_0,T).
\tag{17.9}\label{eq:zeta-functional-equation}
\]
The integer \(\chi_\ell\) and the sign \(\varepsilon_\ell\) are independent
of \(\ell\).
\end{theorem}

\begin{proof}
Apply the abstract determinant computation
\cref{thm:abstract-functional-equation} to the graded spaces
\(V_{i,\ell}\), the operator \(F\), and the perfect pairings
\eqref{eq:poincare-X}, whose multiplier is \(q^d\).  Its alternating
determinant is \(Z(X_0,T)\) by
\cref{thm:cohomological-rationality}.  This proves
\eqref{eq:zeta-functional-equation}; the middle-degree sign is controlled by
the \((-1)^d\)-symmetry of cup product.

If two equations of the displayed shape hold for the same nonzero rational
function \(Z\), their quotient makes
\(T^{\chi_\ell-\chi_{\ell'}}\) constant.  Hence the exponents agree, and then
the constants and signs agree.  Thus both are determined by the
point-counting zeta function and are independent of \(\ell\).
\end{proof}

\section{What remains arithmetic}

\begin{corollary}[Weil theorem before separation of factors]
\label{cor:weil-before-integrality}
The zeta function is rational over \(\mathbf Q\), has the cohomological
factorization \eqref{eq:cohomological-zeta}, has endpoint factors
\eqref{eq:endpoint-factors}, satisfies the functional equation
\eqref{eq:zeta-functional-equation}, and its degree-\(i\) Frobenius
eigenvalues have universal weight \(i\).
\end{corollary}

Three linked assertions remain: every eigenvalue must be shown algebraic over
\(\mathbf Q\), each \(P_{i,\ell}\) must lie in \(\mathbf Z[T]\), and the
polynomial must be independent of \(\ell\).  The next chapter derives all
three from the already proved statements in the safe order
\[
\begin{gathered}
  Z\in\mathbf Z[[T]]\cap\mathbf Q(T)
  \Longrightarrow\text{integral reduced factors}\\
  \Longrightarrow\text{weight separation}\\
  \Longrightarrow P_i\in\mathbf Z[T]\\
  \text{and \(P_i\) is \(\ell\)-independent}.
\end{gathered}
\]

\begin{formalization}
The final assembly has four independent interfaces:
\begin{enumerate}
\item \texttt{perverseProperPurityToCohomology}, implementing
      \[
      \Qbarl[d]\mapsto Ra_*\Qbarl[d]\mapsto
      H^{i-d}(Ra_*\Qbarl[d])=H^i(X);
      \]
\item \texttt{traceFormulaToZeta}, using the determinant--trace identity;
\item \texttt{dualityToFunctionalEquation}, using the similitude pairing;
\item \texttt{universalPurityToSeparatedFactors}, developed in the next
      chapter.
\end{enumerate}
The release theorem must import the proved Deligne--Laumon theorem and proper
purity, not an axiom named “Weil purity.”
\end{formalization}
